\slajdid{09_1-s043}
\begin{frame}[label=09_1-s043]
\gusttekst
Daljim porastom ulaznog napona, napon na drejnu tranzistora T1 opada i kada je
\[V_O=V_{DS,D}=V_{GS,D}-V_{Tn,D}=V_I-V_{Tn,D}\]
i tranzistor $T_D$ ulazi u omsku oblast, pa je zavisnost ulaznog od izlaznog napona
\[\frac{k_{n,L}}2\bigl(2V_{DS,L}(V_{GS,L}-V_{Tn,L})-V_{DS,L}^2\bigr)=\frac{k_{n,D}}2\bigl(2V_{DS,D}(V_{GS,D}-V_{Tn,D})-V_{DS,D}^2\bigr)\]
\[k_{n,L}\bigl(2V_{DS,L}(V_{GS,L}-V_{Tn,L})-V_{DS,L}^2\bigr)=k_{n,D}\bigl(2V_{DS,D}(V_{GS,D}-V_{Tn,D})-V_{DS,D}^2\bigr)\]
uz $V_C=V_{DD}+V_{Tn,L}+\Delta$
\[\alert{1.}\quad k_{n,L}(V_{DD}-V_O)(V_{DD}+2\Delta-V_O)=k_{n,D}V_O(2V_I-2V_{Tn,D}-V_O)\]
\end{frame}
